\chapter{Conclusiones y trabajo futuro}
\label{cap:conclusiones}

\section{Conclusiones}

El desarrollo de la aplicación ha permitido abordar y resolver los problemas de gestión y organización en un hogar compartido. A través de un enfoque centrado en el usuario, se han implementado funcionalidades que facilitan la comunicación, 
la planificación y la gestión de tareas entre los compañeros de piso. La aplicación no solo mejora la eficiencia en la gestión del hogar, sino que también fomenta un ambiente más colaborativo y armonioso entre los inquilinos.

Una vez desarrollada la aplicación y terminado el proceso de pruebas, se ha podido comprobar que la aplicación cumple con los 
objetivos planteados al inicio del proyecto. Los objetivos que se han completado son los siguientes: 

\begin{itemize}
     \item \textbf{Objetivo 1: Análisis de la gestión de hogares compartidos y gestión de gastos:} Este objetivo se cumplió en el capítulo 3 de la memoria, donde se realizó dicho análisis. Gracias a este análisis se pudo identificar mejor los problemas que surgen en los hogares compartidos, y así poder ofrecer unos requisitos más ajustados a la realidad.
     \item \textbf{Objetivo 2: Estudio de aplicaciones similares:} Este objetivo se llevó a cabo también en el capítulo 3, donde se investigó y realizó un análisis de las aplicaciones existentes en el mercado. Esto nos ayudó a saber que funcionalidades eran necesarias para la aplicación y funcionalidades novedosas que podrían ser de interés para los usuarios.
     \item \textbf{Objetivo 3: Análisis de las tecnologías a utilizar:} Este objetivo se realizó en el capítulo 4, donde se realizó un análisis de las diferentes tecnologías (tanto front-end como back-end) que se utilizan en la actualidad. Se realizaron comparativas entre las diferentes tecnologías y se eligieron las más adecuadas para el desarrollo de la aplicación. Esto nos permitió elegir una tecnología con la que estuviéramos seguros y cómodos de trabajar.
     \item \textbf{Objetivo 4: Desarrollo e implementación de la aplicación:} Este fue el objetivo principal del proyecto. Se realizó durante el capítulo 6 y 7 de la memoria. En el capítulo 6 se realizó el análisis y diseño de la aplicación, que esta estrechamente relacionado con el desarrollo de la aplicación, en el, entre otras cosas, se realizaron los requisitos y bocetos de la aplicación, que fue de mucha ayuda para el desarrollo. En el capítulo 7 se realizó el desarrollo de la aplicación, donde se implementaron todas las funcionalidades necesarias para que la aplicación funcionara correctamente.
     \item \textbf{Objetivo 5: Análisis de usabilidad y accesibilidad:} Este objetivo también se realizó en el capítulo 3. Gracias a este análisis se pudo identificar los problemas comunes con la accesibilidad y usabilidad de las aplicaciones, sirvió también para obtener herramientas que nos ayudaron a proporcionar una mejor usabilidad y en un futuro una mejor accesibilidad.
     \item \textbf{Objetivo 6: Validación de la aplicación por usuarios reales:} Este objetivo se completó en el capítulo 8, donde se realizó una evaluación de la aplicación en un entorno real. Se realizaron encuestas a los usuarios para conocer su opinión sobre la aplicación y si cumplía con sus expectativas. Gracias a esto se pudieron identificar posibles mejoras que se podrían realizar en la aplicación y obtener un visión más realista sobre la usabilidad de la aplicación.
\end{itemize}


A parte de los objetivos, se han conseguido logros importantes a lo largo del desarrollo de la aplicación. Estos logros son: 

\begin{itemize}
     \item \textbf{Desarrollo funcional completo:} La aplicación cuenta con todas las funcionalidades básicas necesarias para la gestión de un hogar compartido. Todas las funcionalidades han sido implementadas y probadas, sin obtener errores en su funcionamiento. Lo cual convierte a la aplicación en un producto totalmente funcional.
     \item \textbf{Compatibilidad multiplataforma:} Al haber desarrollado la aplicación en Flutter, la aplicación es compatible con múltiples plataformas tanto Android como iOS, como también versión web y de escritorio. Esto permite que la aplicación sea accesible para una mayor cantidad de usuarios, sin importar el dispositivo que utilicen.
     \item \textbf{Novedad:} La aplicación ofrece funcionalidades novedosas que no se encuentran en otras aplicaciones similares. Esto la convierte en una opción novedosa, única y atractiva para los usuarios. La opción de fusionar tanto la gestión de tareas que se puedan repartir entre los compañeros de piso, como la gestión de gastos, es una funcionalidad que no se encuentra en aplicaciones similares.
     \item \textbf{Usabilidad:} La aplicación ha sido diseñada teniendo en cuenta principalmente la usabilidad. Se han realizado pruebas de usabilidad con usuarios reales, y se han obtenido resultados positivos. La aplicación es fácil de usar y entender, lo que facilita su adopción por parte de los usuarios.
\end{itemize}


\newpage


\section{Trabajos futuros}

Durante todo el desarrollo de este proyecto, se ha asegurado que la funcionalidad de la aplicación fuera la correcta y que no hubiera errores 
en su funcionamiento, por tanto, alguna de las mejoras que se podrían realizar son relacionadas con la accesibilidad de la aplicación y añadir 
alguna de las funcionalidades que nos sugirieron los usuarios durante la evaluación de la aplicación.

A continuación se detallan algunas de las mejoras que se podrían realizar en un futuro:

\begin{itemize}
     \item \textbf{Mejoras en la accesibilidad:} Durante el desarrollo de la aplicación, se tuvo en cuenta la accesibilidad, pero el principal foco fue la usabilidad. Por tanto , se podrían realizar mejoras en la accesibilidad de la aplicación, como por ejemplo, añadir un modo oscuro, o mejorar el contraste de los colores utilizados en la aplicación. También se podría añadir un lector de pantalla para facilitar el uso de la aplicación a personas con discapacidad visual.
     \item \textbf{Nuevas funcionalidades:} Durante la evaluación de la aplicación, los usuarios sugirieron algunas funcionalidades que podrían ser interesantes para añadir a la aplicación. Algunas de estas funcionalidades son: poder añadir el precio de los productos que se comprar y meterlo directamente en el apartado de gastos, podemos clasificar los gastos por categorías, poder asignar las tareas semanales a un día concreto. Estas funcionalidades podrían ser interesantes para añadir a la aplicación en un futuro. 
     \item \textbf{Manual de usuario:} Finalmente, con todas las funcionalidades implementadas, se podría realizar un manual de usuario para la aplicación. Se podría mostrar la primera vez que se inicia en la aplicación y que hubiera una sección permanente en la aplicación donde se pudiera acceder a este manual. Esto ayudaría a los usuarios a entender mejor la aplicación y a utilizarla de forma más eficiente.
     \item \textbf{Lanzamiento de la aplicación:} Finalmente, se podría lanzar la aplicación a diferentes plataformas como Google Play Store o Apple Store. Esto sería el paso final. 
\end{itemize}

Completando todas estas mejoras, la aplicación sería perfecta para usarla en la vida real. De esta forma, se podría dar por finalizado el proyecto. 

